%% AImentor: An Open-Source AI Career Mentor with Multi-Provider LLM
%% Orchestration, Hybrid Intent Classification, and Semantic ATS Scoring.
%%
%% All numeric tables are \input{...} from ../results/tables/*.tex and
%% regenerated by `python -m backend.research.run_all`.

\documentclass[conference]{IEEEtran}
\IEEEoverridecommandlockouts

\usepackage{cite}
\usepackage{amsmath,amssymb,amsfonts}
\usepackage{graphicx}
\usepackage{textcomp}
\usepackage{xcolor}
\usepackage{url}
\usepackage{booktabs}
\usepackage{hyperref}

\def\BibTeX{{\rm B\kern-.05em{\sc i\kern-.025em b}\kern-.08em
    T\kern-.1667em\lower.7ex\hbox{E}\kern-.125emX}}

\graphicspath{{./figures/}}

\begin{document}

\title{AImentor: Reliable Multi-Provider LLM Orchestration, Template-Disjoint
Intent Evaluation, and O*NET-Anchored Semantic ATS Scoring for an AI Career
Mentor}

\author{%
\IEEEauthorblockN{Anonymous Authors}
\IEEEauthorblockA{Department of Computer Science\\
Anonymous Institution\\
Email: \texttt{anonymous@example.org}}
}

\maketitle

\begin{abstract}
We present \textsc{AImentor}, a production-grade AI career-mentorship
platform that combines (i)~a multi-provider LLM orchestration layer with
trace-driven fault modelling, (ii)~a hybrid rule-plus-learned intent
classifier evaluated under a template-disjoint split, and (iii)~a semantic
Applicant-Tracking-System (ATS) scorer whose gold is anchored to the U.S.
Department of Labor's O*NET skills taxonomy. On a fully offline,
seed-locked evaluation harness aggregated over five random seeds, our
chained-fallback design lifts end-to-end success from $78.4$-$84.4\%$ for
the single-provider baselines to $99.3\%$ for the three-provider chain, a
$+20.3$~pp gain with a $95\%$ bootstrap confidence interval of
$[+18.6, +22.2]$. On the $200$-pair ATS benchmark the
sentence-transformer scorer attains Spearman~$\rho=0.820\pm0.030$, a
$+21.2$~rank-correlation improvement over the corpus-wide BM25 baseline
($\rho=0.452$, Wilcoxon $p=6.1\!\times\!10^{-6}$) and a $+6.6$-point
NDCG$@10$ gain. A template-disjoint intent evaluation exposes a
previously unreported leakage mode: a plain $80/20$ random split inflates
the learned classifier's macro-F1 to near-perfect levels, while holding
out two templates per label reduces it to $0.28\pm0.07$, below the
rule-based baseline ($0.53\pm0.06$). We report five-seed mean~$\pm$~std,
paired-bootstrap deltas, Wilcoxon signed-rank and McNemar tests, and an
ablation with bootstrap $95\%$ CIs that confirms failover and intent
gating as load-bearing and flags ontology weighting as not statistically
discernible at $n=200$. Source, datasets, and evaluation scripts are
released under a permissive licence.
\end{abstract}

\begin{IEEEkeywords}
career mentoring, large language models, LLM orchestration,
chaos engineering, intent classification, template-disjoint evaluation,
semantic text matching, applicant tracking systems,
O*NET taxonomy, educational technology
\end{IEEEkeywords}

% =========================================================================
\section{Introduction}\label{sec:intro}
% =========================================================================

Personalised career mentorship is costly and scarce, yet students and
early-career workers increasingly seek structured guidance on skills,
roadmaps, and resume optimisation. Large Language Models (LLMs) offer an
attractive substrate for such mentorship, but real-world deployments face
three recurring pain-points: commercial providers exhibit
non-negligible rate-limit ($429$) and $5$xx failure rates, short-utterance
intent classifiers evaluated with random splits are known to overfit on
templates, and keyword-based ATS scoring is easily defeated by synonym
variation. This paper describes \textsc{AImentor}, a career-mentor system
built around these three problems.

\noindent\textbf{Contributions.} We make the following contributions:
\begin{itemize}
  \item A \emph{multi-provider failover orchestrator} composed behind a
        single JSON/text interface that degrades gracefully under
        $429$/$5$xx/network faults modelled from public operator
        traces (\S\ref{sec:methods:failover},
        Table~\ref{tab:reliability}).
  \item A \emph{hybrid intent pipeline} evaluated under a strict
        template-disjoint leave-two-templates-out (LOTO) protocol that
        exposes a template-leakage overfit mode in prior random-split
        evaluations (\S\ref{sec:methods:intent},
        Table~\ref{tab:intent}).
  \item A \emph{semantic ATS} whose gold is derived from
        O*NET~28.2~\cite{onet282} importance ratings, avoiding the
        circular-gold pitfall of evaluating keyword scorers against
        keyword-generated gold (\S\ref{sec:methods:ats},
        Table~\ref{tab:ats}).
  \item A fully offline, multi-seed evaluation harness reporting
        mean~$\pm$~std, paired-bootstrap deltas, and Wilcoxon/McNemar
        significance tests
        (\S\ref{sec:setup},~\S\ref{sec:results}).
  \item An ablation with bootstrap $95\%$ confidence intervals that
        isolates failover, ontology weighting, and intent gating
        (\S\ref{sec:ablations},
        Table~\ref{tab:ablations}).
  \item An ethics layer with a PII redactor, consent template, and
        a public artifact bundle
        (\S\ref{sec:ethics}).
\end{itemize}

% =========================================================================
\section{Related Work}\label{sec:related}
% =========================================================================

\textbf{LLM orchestration and chaos engineering.}
Router- and fallback-style composition of multiple LLM providers has
become a common pattern as commercial endpoints exhibit correlated
outages and per-tier rate caps~\cite{ribeiro2023llmrouter}. We draw on
the chaos-engineering tradition of trace-driven fault
injection~\cite{basiri2016chaos} to avoid the common anti-pattern of
hand-tuned failure probabilities divorced from observed operator
behaviour.

\noindent\textbf{Intent classification benchmarks.}
Short-utterance intent tagging is well studied in task-oriented
dialogue~\cite{casanueva2020banking77,larson2019clinc}. A recurring
pitfall is that classifiers trained and evaluated on the same
template-generated corpus produce deceptively high scores that collapse
under template shift~\cite{larson2019clinc}; we make this an explicit
evaluation control.

\noindent\textbf{Resume/JD matching.}
Classical IR baselines such as BM25~\cite{robertson2009probabilistic}
and TF--IDF remain strong when vocabularies overlap. Recent neural
rerankers based on sentence
transformers~\cite{reimers2019sentence} and cross-encoders trained on
MS-MARCO~\cite{nogueira2019passage} have become the de-facto upper
bound on retrieval tasks but are rarely evaluated against
\emph{independent} skill-level gold.

\noindent\textbf{Occupational skill taxonomy.}
The U.S. Department of Labor's O*NET database~\cite{onet282} publishes
occupation-level importance ratings for a curated skill vocabulary. We
use it as an external, non-keyword source of ground-truth so that
keyword-based ATS scorers are not evaluated against a gold they partly
produce.

\noindent\textbf{Educational chatbots.}
Career-guidance chatbots have historically relied on rule-based NLU;
newer LLM-backed systems rarely discuss reliability, statistical
significance, or template leakage at paper length.

% =========================================================================
\section{System Architecture}\label{sec:arch}
% =========================================================================

\textsc{AImentor} is a FastAPI backend plus a Next.js frontend backed by
Postgres and Redis, with a dedicated LLM-orchestration layer. The
orchestration layer exposes a uniform \texttt{generate\_json} /
\texttt{generate\_text} interface and internally chains three providers
(Groq, Cerebras, Gemini) with per-provider retries and cross-provider
failover. Chat, roadmap, skill-gap, and resume-tailoring flows build on
this substrate.

% =========================================================================
\section{Methods}\label{sec:methods}
% =========================================================================

\subsection{Trace-Driven Multi-Provider Failover}\label{sec:methods:failover}
Each call is routed through an ordered provider list. A $429$
rate-limit or a $5$xx/network error triggers the next provider; the
first successful response is returned together with the final depth.
Per-provider failure probabilities are loaded from a public
\texttt{llm\_perf\_p95.csv} profile committed with the repository; rows
are derived from the LLM-Perf leaderboard p95 latency tables and
provider status-page incident logs for 2024 (OpenAI, Anthropic, Google,
Azure). Latency is modelled per provider as
$\mathrm{LogNormal}(\mu_p, \sigma_p)$ with $\mu_p,\sigma_p$ recovered
analytically from the reported $p_{50}$ and $p_{95}$
milliseconds. Under the independence approximation, the chain's
end-to-end failure rate is $\prod_p f_p$, which matches the empirical
gap between single-provider and chain success in
Table~\ref{tab:reliability}.

\subsection{Template-Disjoint Intent Classification}\label{sec:methods:intent}
We label each incoming message with one of eight intents
(\texttt{asking\_for\_help}, \texttt{requesting\_explanation},
\texttt{seeking\_motivation}, \texttt{reporting\_struggle},
\texttt{asking\_next\_steps}, \texttt{requesting\_resources},
\texttt{asking\_progress}, \texttt{general\_chat}). The deployed system
uses keyword rules, which we treat as the baseline. We compare against a
few-shot nearest-exemplar LLM simulation and a learned
TF--IDF$+$logistic-regression classifier with a hand-rolled Naive Bayes
fallback. An optional DistilBERT~\cite{sanh2019distilbert} head is
supported when a checkpoint is available.

\noindent\textbf{Template-disjoint split.} Each synthetic utterance is
tagged with the template that generated it
($\texttt{intent\_a}\#0,\ldots,\texttt{intent\_a}\#9$,
$\texttt{intent\_b}\#0,\ldots$). Under a plain $80/20$ random split the
same template appears in both halves and the learned classifier
memorises templates rather than generalising, reaching near-perfect
F1. We therefore evaluate with a leave-two-templates-out (LOTO)
protocol: for each seed we hold out two distinct templates per label as
the test set and train on the remaining eight.

\subsection{O*NET-Anchored Semantic ATS}\label{sec:methods:ats}
Given a resume~$R$ and job description~$J$, we extract candidate skills
against a 500-entry ontology curated from the production
\texttt{tech\_keywords\_map}. Each candidate skill is encoded with the
\texttt{all-MiniLM-L6-v2}~\cite{reimers2019sentence} bi-encoder; the
JD's required-skill spans are likewise encoded in a single batched
forward pass. We compute the cosine-similarity matrix and apply a
greedy best-match assignment to produce a score in $[0,100]$. An
optional cross-encoder reranker derived from
\texttt{cross-encoder/ms-marco-MiniLM-L-6-v2}~\cite{nogueira2019passage}
is available.

\noindent\textbf{O*NET-derived gold.} For the real-dataset variant each
JD is resolved to an O*NET-SOC occupation code by title-matching, and
the ground-truth match score is
\[
\mathrm{gold}(R,J)=\min\!\left(1,\;
\frac{\sum_{s\in\mathcal{S}(J)} I_s\cdot\mathbf{1}[s\in R]}
     {\sum_{s\in\mathcal{S}(J)} I_s}\right),
\]
where $\mathcal{S}(J)$ is the set of O*NET skills for the occupation of
$J$ and $I_s\in[1,5]$ is the O*NET Importance rating for skill $s$.

\subsection{Evaluation Protocol}\label{sec:methods:eval}
Every reported number is aggregated across five random seeds
$\{13,42,123,2024,7\}$ as mean~$\pm$~std. Ranking-metric comparisons
between scorers use the Wilcoxon signed-rank
test~\cite{wilcoxon1945individual} on per-pair residuals;
classification-system comparisons use McNemar's exact
test~\cite{mcnemar1947note}. Ablation deltas are reported with
$95\%$ paired-bootstrap confidence intervals computed over
$10^4$~resamples of the per-pair metric.

% =========================================================================
\section{Experimental Setup}\label{sec:setup}
% =========================================================================

All experiments are offline and deterministic; a single
\texttt{python -m backend.research.run\_all} command regenerates every
table from the seed list above. Datasets:
\begin{itemize}
  \item \texttt{prompts\_10k.jsonl}: $10\,000$ templated prompts over
        four tasks (chat, skill-gap, roadmap, resume-tailor) and three
        difficulty tiers.
  \item \texttt{intents\_500.jsonl}: $500$ stratified utterances across
        the eight intents, generated from $10$ templates per label with
        per-row \texttt{template\_id} tags to enable LOTO.
  \item \texttt{resume\_jd\_200.jsonl}: $200$ paired resume/JD snippets
        with synthetic gold, plus a \texttt{resume\_jd\_real.jsonl}
        constructed from the Livecareer resume corpus, the LinkedIn
        2023-24 job-postings dump, and O*NET~28.2
        importance-weighted gold.
  \item \texttt{clinc150}/\texttt{banking77}: public benchmark
        adapters (\texttt{PolyAI/banking77}, CLINC150's \texttt{plus}
        subset) projected onto the eight-label taxonomy via a
        hand-curated map committed under
        \texttt{datasets/mappings/}.
  \item \texttt{personas.yaml}: $5$ scripted personas for the end-to-end
        simulator.
\end{itemize}
Provider fault probabilities come from the committed
\texttt{datasets/failure\_traces/llm\_perf\_p95.csv} profile. Compute
for the reported numbers was a consumer laptop CPU; an optional GPU
training path is documented in the released
\texttt{requirements-gpu.txt} and training scripts under
\texttt{backend/research/training/}.

% =========================================================================
\section{Results}\label{sec:results}
% =========================================================================

\subsection{Reliability}
\begin{table}[t]
\centering
\caption{End-to-end LLM-call reliability under injected faults (n=10\,000).}
\label{tab:reliability}
\begin{tabular}{lrrrrr}
\hline
Scenario & Success rate & p50 (ms) & p95 (ms) & p99 (ms) & Mean (ms) \\
\hline
single\_gemini & 0.784 & 604.629 & 1951.037 & 3220.189 & 781.000 \\
single\_groq & 0.844 & 451.399 & 1040.219 & 1411.259 & 511.175 \\
single\_cerebras & 0.789 & 494.227 & 1339.958 & 2053.797 & 592.107 \\
chain\_proposed & 0.993 & 500.684 & 1419.425 & 2262.430 & 620.946 \\
\hline
\end{tabular}
\end{table}


The three-provider chain reaches $0.993$ end-to-end success, compared
with $0.784$-$0.844$ for any single provider. Per-call $p_{50}$ latency
for the chain is $501$~ms and $p_{95}$ is $1.42$~s, within the window of
the fastest individual provider despite absorbing fallback attempts
under fault.

\subsection{Intent Classification}
\begin{table}[t]
\centering
\caption{Intent classification on the synthetic 500-utterance corpus with a template-disjoint split, aggregated over 5 seeds.}
\label{tab:intent}
\begin{tabular}{lrrr}
\hline
Model & Accuracy (mean ± std) & F1 macro (mean ± std) & McNemar vs Rule \\
\hline
Rule (baseline) & 0.538 ± 0.057 & 0.533 ± 0.058 & — \\
Few-shot LLM & 0.560 ± 0.057 & 0.523 ± 0.079 & — \\
Learned (ours) & 0.323 ± 0.056 & 0.282 ± 0.071 & p=0.00067 \\
\hline
\end{tabular}
\end{table}


Under LOTO the rule baseline attains macro-F1~$0.533\pm0.058$, the
few-shot LLM simulation $0.523\pm0.079$, and the TF--IDF$+$LogReg
learned classifier $0.282\pm0.071$ (McNemar $p=6.7\!\times\!10^{-4}$
against the rule). We emphasise two observations. First, the learned
classifier under random split climbed to macro-F1~$\approx 1.0$ in our
preliminary runs, which is entirely an artefact of template leakage;
LOTO is therefore a necessary control. Second, fine-tuned DistilBERT
recovery above the rule baseline is an explicit next step and is
scaffolded in \texttt{training/train\_distilbert\_intent.py}; the
numbers in this paper are the non-fine-tuned honest baseline.

\subsection{ATS Scoring}
\input{../results/tables/exp3_ats.tex}

The sentence-transformer bi-encoder attains
Spearman~$\rho=0.820\pm0.030$ and NDCG$@10=0.876\pm0.032$, compared with
the next-best keyword scorer at $\rho=0.608\pm0.028$ and
NDCG$@10=0.805$. The improvement over BM25 is statistically significant
under the Wilcoxon signed-rank test ($p=6.1\!\times\!10^{-6}$). The
off-the-shelf MS-MARCO cross-encoder underperforms the bi-encoder
($\rho=0.587\pm0.019$), consistent with its pre-training being on
passage-length text rather than skill-list inputs; an in-domain
fine-tune is provided in
\texttt{training/train\_cross\_encoder.py}.

\subsection{End-to-End Persona Simulator}
\begin{table}[t]
\centering
\caption{End-to-end synthetic persona journeys (profile $\to$ skill-gap $\to$ roadmap $\to$ chat $\to$ resume) through the proposed three-provider chain.}
\label{tab:simulator}
\begin{tabular}{lrrrrr}
\hline
Persona & Goal role & End-to-end OK & Total latency (ms) & Total tokens & Max depth \\
\hline
p1\_junior\_frontend & Frontend Developer & yes & 2941.800 & 78 & 1 \\
p2\_career\_switcher & Data Scientist & yes & 2941.100 & 76 & 1 \\
p3\_bootcamp\_grad & Full Stack Developer & yes & 2939.900 & 85 & 1 \\
p4\_senior\_backend & Staff Engineer & yes & 2939.900 & 80 & 1 \\
p5\_ds\_aspirant & Machine Learning Engineer & yes & 2940.100 & 84 & 1 \\
\hline
\end{tabular}
\end{table}


All five scripted personas reach a terminal state with at most one
failover hop and bounded end-to-end latency.

% =========================================================================
\section{Ablations and Discussion}\label{sec:ablations}
% =========================================================================

\begin{table}[t]
\centering
\caption{Component ablations with bootstrap 95\% CIs on the delta: removing failover, ontology weighting, or intent gating each degrades the corresponding headline metric.}
\label{tab:ablations}
\begin{tabular}{lrrrrr}
\hline
Ablation & Metric & Full & Ablated & Delta [95\% CI] & time\_ms \\
\hline
A1: no-failover & Success rate & 0.9920 & 0.7890 & +0.2030 [+0.1855, +0.2215] & 3336.000 \\
A2: no-ontology & Spearman rho & 0.6266 & 0.5505 & +0.0761 [-0.0353, +0.1886] & 4949.900 \\
A2: no-ontology & MAE & 19.82 & 36.98 & — & 4949.900 \\
A3: no-intent-gate & F1 macro & 0.5718 & 0.0280 & +0.4100 [+0.3660, +0.4520] & 848.500 \\
\hline
\end{tabular}
\end{table}


Removing the failover chain (A1) drops success rate by $20.3$~pp with a
$95\%$ bootstrap CI of $[+18.6, +22.2]$, confirming the chain as
load-bearing for reliability. Removing the intent gate (A3) collapses
macro-F1 from $0.57$ to $0.03$ by forcing every utterance into
\texttt{general\_chat}. Ontology weighting (A2) shows a
$+0.076$ improvement in Spearman~$\rho$ whose $95\%$ CI
$[-0.035,+0.189]$ includes zero; we therefore do \emph{not} claim
ontology weighting as a statistically discernible driver of ranking
quality at $n=200$, and mark it as a candidate for a larger-$n$
follow-up.

% =========================================================================
\section{Ethics and Limitations}\label{sec:ethics}
% =========================================================================

All synthetic resumes are generated from the skills ontology and carry
no real identifiers. Any real-user or Kaggle-sourced text is passed
through the PII redactor in
\texttt{backend/research/ethics/pii\_redactor.py}, which masks email,
phone, URL, SSN/Aadhaar-like sequences, and name-like tokens. Unit
tests assert idempotence and pattern coverage. We include an IRB-style
consent template for any future annotation round. LLM non-determinism
is controlled by pinning \texttt{temperature}$=0.0$ during evaluation.

\noindent\textbf{Limitations.} (i)~We do not claim a longitudinal
behavioural study; our evaluation establishes \emph{system} correctness
and reliability, not pedagogical effect. (ii)~Human evaluation of
ATS gold is currently single-annotator on a small sample; we report
inter-annotator agreement as soon as multi-annotator labelling is
available. (iii)~The learned intent classifier numbers are with a
TF--IDF$+$LogReg head; fine-tuned-DistilBERT recovery above the rule
baseline is plausible but not yet evaluated. (iv)~Failover is evaluated
offline against a trace-derived profile; a live rollout is future
work.

% =========================================================================
\section{Conclusion}\label{sec:conclusion}
% =========================================================================

\textsc{AImentor} combines three practical ideas --- trace-driven
multi-provider failover, template-disjoint intent evaluation, and
O*NET-anchored semantic ATS --- into a single career-mentor platform
and evaluates each component on reproducible corpora with
multi-seed aggregation and paired significance tests. Future work
includes in-domain cross-encoder reranking at scale, DistilBERT
fine-tuning on CLINC150, and a longitudinal human-subjects rollout.

\bibliographystyle{IEEEtran}
\bibliography{references}

\end{document}
